% Short report on the black-hole and neutron-star production runs of 2026-09-17/18.
%
% Build:  cd Report/populations && latexmk -pdf populations_report.tex
% Figures are pulled from ../../figures/prod_20260918/ via \graphicspath below; they are the
% exact files the analysis scripts wrote, not copies, so this document cannot drift from them.
\documentclass[english,10pt,a4paper]{article}
\usepackage[T1]{fontenc}
\usepackage[margin=2.4cm]{geometry}
\usepackage{graphicx}
\graphicspath{{./}{../../figures/prod_20260918/}{../../figures/yield_prefix_20260922/y3/}}
\usepackage{mathtools}
\usepackage{amssymb}
\usepackage{booktabs}
\usepackage{xcolor}
\usepackage[colorlinks=true,allcolors=blue]{hyperref}
\usepackage{caption}
\usepackage[authoryear,round]{natbib}
\usepackage[section]{placeins}

\newcommand{\tE}{t_{\rm E}}
\newcommand{\piE}{\pi_{\rm E}}
\newcommand{\tetE}{\theta_{\rm E}}
\newcommand{\Ml}{M_{\rm L}}
\newcommand{\Msun}{M_\odot}
\newcommand{\Neff}{N_{\rm eff}}

\title{Compact objects in the Roman--Rubin joint microlensing forecast:\\
       black holes and neutron stars}
\author{Ali Salehi}
\date{\today}

\begin{document}
\maketitle

\begin{abstract}
We report two production runs of the joint Roman\,+\,Rubin Galactic-bulge microlensing
simulation, one for a population of isolated black holes with a log-uniform mass function over
$3$--$1000\,\Msun$, and one for neutron stars drawn from a Gaussian at $1.35\pm0.15\,\Msun$
truncated to $1.10$--$2.20\,\Msun$. Each run scans 1829 sightlines over $\sim\!68\,\mathrm{deg}^2$
and characterises every detected event three times -- with Rubin alone, with Roman alone, and
with the two combined -- through independent Fisher matrices. All pooled quantities are
event-rate weighted and quoted with the \citet{Kish1965} effective sample size. We find that the joint fit
improves the timescale by a factor $6$--$9$ over Rubin alone and $\sim\!1.35$ over Roman
alone, and that it barely improves $\tetE$ at all, which is set by astrometry and therefore by
Roman. The centroid shift exceeds Roman's per-exposure astrometric floor for $90\%$ of
black-hole events and $0.6\%$ of neutron-star events. Satellite parallax from L2 is found to be
almost worthless for both populations. Following \citet{SajadianMakler2026}, we compute the
probability that the two lensing-induced images are separately resolvable, and find it to be
dominated by the assumed resolution criterion rather than by the lens population. Attaching the
event rate and an assumed abundance $F$ to each draw, we also convert the Monte Carlo sample sizes
into absolute yields: for $F_{\rm BH}=0.01$ Roman sees $\sim\!49$ black-hole events over a decade,
of which $\sim\!4$ have their mass measured to $10\%$. Because the yield is exactly linear in $F$, we also give the yield per unit
abundance for every selection, and the recipe that produces it, so that any $F$ can be used. The normalisation is validated against the
code's own optical depth to $0.1\%$, against OGLE-IV, and against \citet{Penny2019}. Matched to
the mass function and abundance of \citet{SajadianSahu2023}, we detect $\sim\!2.3\times$ more
black holes per unit area and season than they do --- consistent with the final survey design
filling the $2.3$-year season gap their forecast was built around --- but characterise a smaller
fraction of them, a deficit concentrated on the parallax rather than on the astrometry.
\end{abstract}

\section{Simulation and runs}

The simulator draws source stars and lenses along each sightline from Besan\c{c}on \citep{Robin2003}/CMD
population models with three-dimensional dust extinction, generates per-band light curves
against the real Rubin visit list and the published Roman GBTDS cadence, applies a Monte Carlo
detection test, and forecasts parameter precision with a Fisher matrix. The photometric matrix
has nine parameters $(u_0,\tE,f_{b,0},\piE,\xi,t_0,m_{b,0},f_{b,1},m_{b,1})$ and the astrometric
matrix four $(\tetE,\mu_{s,1},\mu_{s,2},\piE)$; which parameters are free is decided per event
from the epochs each survey actually contributed. Crucially there is one such pair of matrices
\emph{per survey partition} -- Rubin-only, Roman-only, and joint -- so that ``what the
combination adds'' is measured rather than assumed.

Both runs used identical configuration apart from the lens mass function
(\verb|--events 300 --lenses 50 --stride-roman 5 --pair-satellite|), a fixed detection
threshold $\Delta\chi^2 \ge 500$, and the same RNG seed. The stratified grid places sightlines
every $0.1^\circ$ inside Roman's footprint and every $0.2^\circ$ outside it, which oversamples
the footprint by design; the per-sightline sky area $w_{\rm area}$ undoes that in every pooled
statistic. Table~\ref{tab:runs} summarises the outcome.

\begin{table}[h]
\centering
\caption{Run configuration and outcome. ``Barren'' sightlines drew stars but characterised
nothing; they run to the full draw cap and so dominate the raw row count while contributing no
detections whatsoever. The weightable sample is the one every number in this report is computed
from. ``Detected draws'' is a \emph{Monte Carlo sample size}, not a forecast: every lens is drawn
from the population under test, i.e.\ $F=1$ implicitly. See Section~\ref{sec:yield} for the
conversion to an absolute yield.}
\label{tab:runs}
\begin{tabular}{lrr}
\toprule
 & black holes & neutron stars \\
\midrule
mass function            & log-uniform          & Gaussian \citep{OzelFreire2016} \\
mass range [$\Msun$]     & $3$--$1000$          & $1.10$--$2.20$ \\
sightlines scanned       & 1829                 & 1829 \\
\quad aggregated         & 1612                 & 1612 \\
\quad no coverage        & 140                  & 140 \\
\quad barren             & 77                   & 77 \\
\quad hit draw cap       & 77                   & 77 \\
rows written             & 4\,679\,203          & 5\,012\,349 \\
weightable draws         & 829\,202             & 1\,162\,348 \\
detected draws (MC)      & 110\,144             & 93\,685 \\
$\Neff$ (weighted)       & 143\,499             & 296\,562 \\
CPU time                 & $\sim\!13.5$\,h      & $\sim\!13.5$\,h \\
\bottomrule
\end{tabular}
\end{table}

\paragraph{Weighting.} The Monte Carlo draws the lens \emph{distance} in proportion to the event
rate but the lens \emph{mass} from the number density and the velocities from plain Gaussians,
so it is missing the rate's $\sqrt{\Ml}\,v_t$ factor. Every pooled fraction, median or ratio in
this report therefore carries the per-event weight
$W \propto w_{\rm area}\,N_\star/n_{\rm sim}\,\sqrt{\Ml}\,v_t\,Z(D_s)$, and every weighted number
is quoted with $\Neff = (\sum w)^2/\sum w^2$, because the weight costs precision. Validation:
weighting all draws of an earlier bulge run reproduces the OGLE-IV efficiency-corrected mean
timescale of $22$\,d (we obtain $24.0$\,d; unweighted the same sample gives $56.2$\,d).

\section{What each survey contributes}

\begin{figure}[h]
\centering
\includegraphics[width=\textwidth]{p6_synergy.png}
\caption{(a) Weighted share of detected events found by each survey. (b) Median ratio of the
joint to the single-survey timescale error, \emph{restricted to events that both surveys
characterised independently}. That restriction is essential: Roman contributes epochs to only a
few per cent of draws, so pooled over all characterised events the ``joint'' fit is literally the
Rubin fit for the great majority and the ratio is $1.000$ by construction.}
\label{fig:synergy}
\end{figure}

Rubin finds the overwhelming majority of events in both populations -- it covers the whole
$68\,\mathrm{deg}^2$ scan, while Roman's footprint is $\sim\!2\%$ of it -- but the events the two
share are the ones that are well measured. Table~\ref{tab:synergy} gives the paired comparison.

\begin{table}[h]
\centering
\caption{Detection provenance and the joint-fit gain. The $\sigma$ ratios are weighted medians
over events characterised independently by \emph{both} surveys (15\,849 black-hole and 12\,569
neutron-star events); values below unity mean the joint fit is better.}
\label{tab:synergy}
\begin{tabular}{lrrrr}
\toprule
 & \multicolumn{2}{c}{black holes} & \multicolumn{2}{c}{neutron stars} \\
\cmidrule(lr){2-3}\cmidrule(lr){4-5}
 & vs.\ Rubin & vs.\ Roman & vs.\ Rubin & vs.\ Roman \\
\midrule
detected by Rubin only [\%] & \multicolumn{2}{c}{89.4} & \multicolumn{2}{c}{90.0} \\
detected by Roman only [\%] & \multicolumn{2}{c}{5.3}  & \multicolumn{2}{c}{6.3} \\
detected by both [\%]       & \multicolumn{2}{c}{5.2}  & \multicolumn{2}{c}{3.6} \\
\midrule
median $\sigma(\tE)$ ratio    & 0.165 & 0.733 & 0.116 & 0.742 \\
median $\sigma(\piE)$ ratio   & 0.108 & 0.826 & 0.114 & 0.798 \\
median $\sigma(\tetE)$ ratio  & 0.187 & 0.979 & 0.152 & 0.975 \\
\bottomrule
\end{tabular}
\end{table}

The asymmetry is the result. Roman's dense cadence sharpens what Rubin alone could do on the
timescale by a factor $6$--$9$ and on the parallax by $\sim\!9$. Rubin's decade-long baseline
sharpens Roman by a more modest $\sim\!1.25$--$1.35$. But for $\tetE$ the ratio against Roman
alone is $0.975$--$0.979$: Rubin adds essentially nothing, because $\tetE$ is an astrometric
quantity and Roman is the only astrometric instrument in this pairing. Any science case resting
on $\tetE$ -- which includes the lens mass, through $\Ml = \tetE/(\kappa\piE)$ -- is a Roman
case, not a joint one.

\section{Fisher forecast precision}

\begin{figure}[h]
\centering
\includegraphics[width=\textwidth]{p7_precision.png}
\caption{Weighted cumulative distribution of the fractional parameter error, per survey and
population. \textbf{Each curve is over its own characterised set}, not a common one: Roman
characterises only inside its footprint, where coverage is good, while the joint curve also
contains every Rubin-only event. That is why Roman's $\tetE$ curve lies above the joint curve in
panel (c), which is a difference of samples and not a violation of
$\sigma_{\rm joint}\le\sigma_{\rm single}$.}
\label{fig:precision}
\end{figure}

\begin{table}[h]
\centering
\caption{Weighted fraction of characterised events measured to better than $10\%$. Read down a
column, not across: each entry has its own denominator (see Fig.~\ref{fig:precision}).}
\label{tab:tenpercent}
\begin{tabular}{lrrrrrr}
\toprule
 & \multicolumn{3}{c}{black holes [\%]} & \multicolumn{3}{c}{neutron stars [\%]} \\
\cmidrule(lr){2-4}\cmidrule(lr){5-7}
 & joint & Rubin & Roman & joint & Rubin & Roman \\
\midrule
$\tE$   & 19.77 & 17.66 & 18.61 & 7.71 & 6.12 & 10.51 \\
$\piE$  &  3.13 &  2.14 &  7.71 & 1.50 & 1.08 &  2.91 \\
$\tetE$ & 43.87 & 41.11 & 81.41 & 6.03 & 0.37 & 42.28 \\
$\Ml$   &  1.50 &  0.44 &  7.07 & 0.34 & 0.01 &  1.85 \\
\bottomrule
\end{tabular}
\end{table}

The lens mass is the hardest quantity in the table, as expected: it needs both $\tetE$ and
$\piE$, and $\piE$ is the weaker of the two. Within Roman's footprint, where the paired
comparison is meaningful, the median $\sigma(\Ml)/\Ml$ improves by a factor $0.121$ (black
holes) and $0.138$ (neutron stars) relative to Rubin alone, but only to $0.862$ and $0.806$
relative to Roman alone -- the same asymmetry as $\tetE$, and for the same reason. Note the
footprint sample is much smaller than the full scan: $\Neff = 5471$ and $5372$ respectively, so
these are the least precise numbers in this report.

\begin{figure}[h]
\centering
\includegraphics[width=\textwidth]{p7_precision_te.png}
\caption{Median fractional error against the event timescale. The two surveys fail at opposite
ends: Rubin's $\sim\!3$-day cadence under-resolves short events, while Roman's $72$-day seasons
cannot contain long ones, which is where Rubin's decade of coverage matters.}
\label{fig:precision_te}
\end{figure}

\section{Astrometric microlensing}

\begin{figure}[h]
\centering
\includegraphics[width=\textwidth]{p6_astrometry.png}
\caption{(a) Weighted distribution of the maximum centroid shift reached by each event. (b) The
blend-diluted shift against lens mass, with the weighted interquartile band; the dotted line is
Roman's $1.1$\,mas per-exposure astrometric floor.}
\label{fig:astrometry}
\end{figure}

The centroid of the unresolved image pair is displaced by
$\delta\theta(u) = \tetE\,u/(u^2+2)$, which peaks at $u=\sqrt{2}$ and \emph{not} at closest
approach; an event whose trajectory crosses $u=\sqrt{2}$ therefore reaches $\tetE/\sqrt{8}$.
Because $\tetE \propto \sqrt{\Ml}$, this is the channel where the two populations diverge most
sharply (Table~\ref{tab:ast}): black holes cross Roman's per-exposure floor at
$\Ml \approx 20\,\Msun$ and $90\%$ of their detections exceed it, while neutron stars sit an
order of magnitude below it. Blending is a small correction here -- the Roman-band source flux
fraction averages $0.95$ -- but it is included rather than assumed negligible.

\begin{table}[h]
\centering
\caption{Astrometric shift, satellite parallax, and image resolvability. The resolution
probability is the weighted fraction of that survey's detections having at least three epochs at
which both images are separately detectable and separated by more than the stated bar.}
\label{tab:ast}
\begin{tabular}{lrr}
\toprule
 & black holes & neutron stars \\
\midrule
median max.\ centroid shift [mas]      & 3.61  & 0.265 \\
95th percentile [mas]                  & 9.75  & 0.610 \\
fraction above $1.1$\,mas floor [\%]   & 90.2  & 0.57 \\
\midrule
satellite parallax, median gain        & 1.000 & 1.000 \\
\quad fraction gaining $>1.1\times$ [\%] & 0.003 & 0.200 \\
\quad fraction gaining $>2\times$ [\%]   & 0.000 & 0.064 \\
\midrule
$P(\text{resolve})$ Rubin, $\mathcal{D}=5$ [\%]  & 16.71  & 0.0015 \\
$P(\text{resolve})$ Roman, $\mathcal{D}=5$ [\%]  & 41.63  & 6.13 \\
$P(\text{resolve})$ Roman, $\mathcal{D}=20$ [\%] & 14.02  & 0.171 \\
$P(\text{resolve})$ Roman, PSF FWHM [\%]         &  0.875 & 0.0001 \\
\bottomrule
\end{tabular}
\end{table}

\section{Satellite parallax}

\begin{figure}[h]
\centering
\includegraphics[width=\textwidth]{p6_parallax.png}
\caption{Survival function of the gain from placing Roman at L2 rather than at Earth, for the
parallax (a) and the lens mass (b). The same event is characterised twice, differing only in the
observer's position, so this is a genuinely controlled comparison. Note the logarithmic vertical
axis: the median gain is exactly $1$, and the entire effect lives in a tail of order $0.1\%$.}
\label{fig:parallax}
\end{figure}

This is a null result and worth stating plainly. The $0.01$\,au displacement of L2 gives a
parallax baseline that is negligible against the Earth's own annual motion for events lasting
months to years, which is what both of these populations produce. Fewer than $0.2\%$ of events
gain even a $10\%$ improvement in $\sigma(\piE)$. Satellite parallax is a technique for
\emph{short} events, and these are not short.

\section{Resolving the two images}

\begin{figure}[h]
\centering
\includegraphics[width=\textwidth]{p6_resolution.png}
\caption{(a) Probability that the two lensing-induced images are resolvable, for three
resolution criteria. (b) The same probability at the most permissive criterion, against lens
mass.}
\label{fig:resolution}
\end{figure}

A point lens forms two images separated by $\Delta\theta(u) = \tetE\sqrt{u^2+4}$. Following
\citet{SajadianMakler2026} we count \emph{epochs} at which both images are individually above the
single-visit depth and their separation exceeds a stated bar, and call the pair resolvable when
at least three epochs qualify. Counting epochs rather than evaluating the separation once per
event is not a refinement but a necessity: the separation is \emph{smallest} at closest approach
and grows as the source departs, while the minor image's magnification
$A_- = (u^2+2)/(2u\sqrt{u^2+4}) - 1/2$ collapses on the same motion, so the images are least
separated exactly when both are bright.

The result is dominated by the choice of criterion. With $\sigma_r = \mathcal{D}\sigma_a$ and
$\mathcal{D}=5$ -- the faint-limit value -- Roman resolves $41.6\%$ of its black-hole detections;
at $\mathcal{D}=20$ this falls to $14.0\%$, and at the diffraction-limited PSF FWHM to $0.87\%$.
Any single number here is a choice of convention presented as a result, which is why all three
are reported. Roman outperforms Rubin by roughly $2.5\times$ at a common $\mathcal{D}$ and by
three orders of magnitude at the PSF bar, reflecting its $105$\,mas F146 PSF against Rubin's
arcsecond seeing disc.

\section{Where the measurable lenses are}

\begin{figure}[h]
\centering
\includegraphics[width=\textwidth]{p7_mass_distance.png}
\caption{Weighted fraction of detected events whose lens mass is measured to better than $10\%$,
in the mass--distance plane. Measurability concentrates on \emph{nearby, massive} lenses, as
$\tetE = \sqrt{\kappa \Ml \pi_{\rm rel}}$ requires.}
\label{fig:massdist}
\end{figure}

A mass precision quoted without saying which lenses it applies to is a statement about a sample
rather than about the Galaxy. Across all detections the mass is measured to $10\%$ for $1.50\%$
of black-hole events (median lens distance $5.15$\,kpc) and $0.34\%$ of neutron-star events
(median $3.66$\,kpc), and Fig.~\ref{fig:massdist} shows those events are drawn overwhelmingly
from lenses within a few kpc.

\begin{figure}[h]
\centering
\includegraphics[width=\textwidth]{p7_sky.png}
\caption{(a) Relative weighted event yield across the scanned region, for the black-hole run.
(b) Weighted fraction of detections with $\sigma(\tE)/\tE < 0.1$ from the joint fit. The yield
concentrates toward the inner bulge; the precision is set by Rubin's visit density and is
comparatively flat.}
\label{fig:sky}
\end{figure}

\section{Absolute event yields}
\label{sec:yield}

Every count quoted so far is a \emph{Monte Carlo sample size}. The simulator draws a fixed number
of lenses per sightline, draws every one of them from the population under test, and reports how
many passed the detection cut. The $110\,144$ and $93\,685$ of Table~\ref{tab:runs} are therefore
the size of the detected sample, not a prediction: they describe a Galaxy in which \emph{all} of
the stellar mass is black holes, or all of it neutron stars. Converting them into a number of
events requires two ingredients the run does not contain --- the event rate carried by each draw,
and the true abundance of the population.

\paragraph{The rate attached to a draw.} For a single lens of mass $\Ml$ at distance $D_l$ moving
with transverse speed $v_t$ past a source at $D_s$, the differential event rate per source star is
$\mathrm{d}\Gamma = 2\theta_{\rm E} D_l\, v_t\, n(D_l)\,\mathrm{d}D_l$ with
$\theta_{\rm E} = \sqrt{\kappa \Ml \pi_{\rm rel}}$ and $\kappa = \sqrt{4G/c^2}$. Because the
simulator already draws $D_l$ in proportion to the rate, what remains is a per-draw factor

\begin{equation}
\gamma_i \;=\; 2u_{0,\max}\,\kappa\,\frac{F}{\langle \Ml\rangle}\;
             Z(D_{s,i})\,\sqrt{M_{{\rm L},i}}\;v_{t,i},
\label{eq:gamma}
\end{equation}

where $Z(D_s)$ normalises the lens column along that sightline, $u_{0,\max}=3$ is the impact
parameter the run samples to, and $F$ is the fraction of the Galaxy's \emph{stellar mass} locked
in the compact-object population. Summing over draws and restoring the sky area and star counts
gives the expected number of events peaking in an observing window of length $T$:

\begin{equation}
N \;=\; T\,\cdot\,2u_{0,\max}\,\kappa\,\frac{F}{\langle \Ml\rangle}
        \sum_i W_i\,\mathrm{det}_i ,
\qquad
W_i \;\propto\; w_{{\rm area},k}\,\frac{N_{\star,k}}{n_{{\rm sim},k}}\,
                \sqrt{M_{{\rm L},i}}\;v_{t,i}\,Z(D_{s,i}),
\label{eq:yield}
\end{equation}

with $\mathrm{det}_i\in\{0,1\}$ the detection flag for the selection of interest and $W_i$ the same
weight already used for every pooled statistic in this report; this section adds only the
normalising constants, which cancel in a ratio and so were never tracked before. We take $T = 10\,\mathrm{yr} - 4\,\mathrm{d}$,
the window over which $t_0$ is drawn.

\paragraph{Why $F$ alone does not set the count.} Equation~(\ref{eq:yield}) depends on
$F/\langle \Ml\rangle$ and on $\sqrt{\Ml}$ inside the weight, never on $F$ by itself. A mass
fraction buys \emph{mass}, but lensing counts \emph{objects}, and a cross-section that grows only
as $\sqrt{\Ml}$. The relevant figure of merit is $\langle\sqrt{\Ml}\rangle/\langle \Ml\rangle$,
which is $1.97$ for ordinary bulge stars ($\langle \Ml\rangle = 0.23\,\Msun$) but only $0.060$ for
the log-uniform $3$--$1000\,\Msun$ black-hole function ($\langle \Ml\rangle = 172\,\Msun$): a
factor $33$ suppression at equal mass fraction, of which detection efficiency returns about
$4\times$ because black-hole events last far longer. This single ratio is why a population holding
a few per cent of the stellar mass contributes a few tenths of a per cent of the events.

\paragraph{Two counting conventions.} The simulator keeps a drawn source with probability equal to
its blend fraction, which counts events \emph{per resolved object} --- the legacy convention of
\citet{SajadianMakler2026} and of \citet{SajadianSahu2023}. Re-weighting by $1/P(\text{kept})$ instead
counts events on \emph{every} star, blended or not. Both are reported below; they differ by
$3$--$8\%$ here, because most detected sources are only mildly blended. Comparisons with the
literature use the per-object convention.

\paragraph{The abundance $F$.} No single value is established, so a grid is quoted. \citet{SajadianMakler2026}
adopt $F_{\rm BH} = 4$--$5\times10^{-3}$. Population syntheses are an order of magnitude
above that: \citet{Olejak2020} give $1.2\times10^{8}$ black holes at $\sim\!14\,\Msun$ and
\citet{Lam2020} $2\times10^{8}$ at $5$--$16\,\Msun$, both implying $F_{\rm BH}\approx0.03$;
\citet{Sweeney2022} place neutron stars and black holes together at $\sim\!1\%$ of the stellar
mass. The bulge census of \citet{Gould2000}, with a mass ratio MS:WD:NS:BH of $69\!:\!22\!:\!6\!:\!3$,
gives $F_{\rm BH}=0.03$ and $F_{\rm NS}=0.06$. We therefore tabulate
$F_{\rm BH}\in\{0.005,0.01,0.03\}$ and $F_{\rm NS}\in\{0.005,0.01,0.03,0.06\}$.

\subsection{Validation of the normalisation}

The rate constants in Eq.~(\ref{eq:gamma}) are checked three ways before being used
(Table~\ref{tab:yieldcheck}). The first is internal and the strongest: the optical depth rebuilt
from these constants, $\tau = (\pi/2u_{0,\max})\langle\gamma \tE\rangle$, is compared against the
optical depth the C++ computes independently for every sightline, and agrees to $0.1\%$ in the
pooled sum. The second is external: the event rate and optical depth for OGLE-like sources
($I<21$, $u_0<1$) are compared against \citet{Mroz2019}. The third fixes the Roman
normalisation against \citet{Penny2019}.

\begin{table}[h]
\centering
\caption{Validation of the absolute normalisation. The OGLE comparison is against the southern
fields of \citet{Mroz2019}, $\Gamma = 14.3\times10^{-6}\exp[0.52(3-|b|)]$ per star per year.
The Roman comparison is per unit area and per day of observing coverage, because Penny's survey
runs $6\times72 = 432$ days while ours runs $10$ seasons totalling $693$ days (counted directly
from the visit list by clustering epochs on gaps $>20$\,d).}
\label{tab:yieldcheck}
\begin{tabular}{llr}
\toprule
check & quantity & result \\
\midrule
internal, $\tau$   & rebuilt $/$ C++ \texttt{opt\_1e6}, bulge run & 0.9985 \\
                   & \quad black-hole run                        & 0.9999 \\
                   & \quad neutron-star run                      & 0.9998 \\
\midrule
OGLE-IV, $\Gamma$  & model $/$ observed at $b = -1.5^\circ$      & 0.97 \\
                   & model $/$ observed at $b = -5.1^\circ$      & 1.47 \\
\midrule
\citet{Penny2019} & events deg$^{-2}$ d$^{-1}$, $|u_0|<1$       & 31.9 \\
                   & \textbf{this work}, $u_0<3$, $\Delta\chi^2\ge500$ & \textbf{40.4} \\
                   & events deg$^{-2}$ d$^{-1}$, $|u_0|<3$       & 63.8 \\
\bottomrule
\end{tabular}
\end{table}

The Penny bracket is predicted rather than fitted, and is the reason it is persuasive: sampling
$u_0<3$ obliges us to exceed the $|u_0|<1$ count, while the $\Delta\chi^2\ge500$ cut rejects most
low-magnification events --- a $u_0=3$ light curve peaks at $A=1.017$ --- obliging us to fall
below the $|u_0|<3$ count. We land between the two.

The OGLE row, by contrast, is a real discrepancy and is reported as one. The model reproduces the
observed rate near the minor axis but overproduces events by $\sim\!40\%$ at $|b|=5^\circ$: its
latitude gradient is too shallow. Yields quoted for the full $59\,\mathrm{deg}^2$ scan inherit
that error and should be read as accurate to a few tens of per cent; yields inside Roman's
footprint, which lies at small $|b|$, are much less affected.

\subsection{Yields}

Tables~\ref{tab:yieldbh} and~\ref{tab:yieldns} give the expected number of events whose peak falls
in the ten-year window, for each survey separately and for the combination. Quoted errors are
Monte Carlo only and do not include the systematic budget discussed above.

\begin{table}[h]
\centering
\caption{Black-hole yields over the ten-year window, per-object convention. The upper block is the
$3$--$1000\,\Msun$ population actually simulated. The lower block re-weights the same draws to a
log-uniform $3$--$50\,\Msun$ function --- \emph{not} the population run, but the range adopted by
\citet{SajadianSahu2023}, included so that the comparison is like for like. ``Roman'' rows are
necessarily confined to the $1.47\,\mathrm{deg}^2$ footprint.}
\label{tab:yieldbh}
\begin{tabular}{lrrr}
\toprule
 & $F=0.005$ & $F=0.01$ & $F=0.03$ \\
\midrule
\multicolumn{4}{l}{\emph{$3$--$1000\,\Msun$, as simulated}}\\
Rubin, full $59\,\mathrm{deg}^2$ scan  & 221   & 441   & 1320 \\
Rubin, within Roman footprint          & 14.1  & 28.2  & 84.6 \\
Roman                                  & 24.6  & 49.2  & 147  \\
both surveys                           & 12.1  & 24.3  & 72.8 \\
Roman, $\sigma(\Ml)/\Ml<10\%$          & 1.89  & 3.77  & 11.3 \\
Roman, $\sigma(\Ml)/\Ml<5\%$           & 0.99  & 1.98  & 5.95 \\
\midrule
\multicolumn{4}{l}{\emph{re-weighted to $3$--$50\,\Msun$, for comparison only}}\\
Rubin, within Roman footprint          & 43.2  & 86.4  & 259  \\
Roman                                  & 64.4  & 129   & 386  \\
Roman, $\sigma(\Ml)/\Ml<10\%$          & 3.97  & 7.94  & 23.8 \\
Roman, $\sigma(\Ml)/\Ml<5\%$           & 1.95  & 3.90  & 11.7 \\
\bottomrule
\end{tabular}
\end{table}

\begin{table}[h]
\centering
\caption{Neutron-star yields over the ten-year window, per-object convention.}
\label{tab:yieldns}
\begin{tabular}{lrrrr}
\toprule
 & $F=0.005$ & $F=0.01$ & $F=0.03$ & $F=0.06$ \\
\midrule
Rubin, full $59\,\mathrm{deg}^2$ scan & 1340 & 2680 & 8040 & 16100 \\
Rubin, within Roman footprint         & 98.3 & 197  & 590  & 1180  \\
Roman                                 & 141  & 281  & 844  & 1690  \\
both surveys                          & 50.8 & 102  & 305  & 610   \\
Roman, $\sigma(\Ml)/\Ml<10\%$         & 3.54 & 7.08 & 21.2 & 42.5  \\
Roman, $\sigma(\Ml)/\Ml<5\%$          & 1.56 & 3.12 & 9.35 & 18.7  \\
\bottomrule
\end{tabular}
\end{table}

\subsection{The yield at any abundance}
\label{sec:yieldF}

Tables~\ref{tab:yieldbh} and~\ref{tab:yieldns} are samples of a single straight line. $F$ enters
Eq.~(\ref{eq:yield}) only through the prefactor $F/\langle\Ml\rangle$, and $\langle\Ml\rangle$ is
fixed by the mass function rather than by how much of it there is, so every selection has
\emph{one} coefficient, its yield at $F=1$, and
\begin{equation}
N(F) \;=\; F\,N_1 .
\label{eq:linear}
\end{equation}
$N_1$ is the Monte Carlo count of Table~\ref{tab:runs} turned into a number of events. It is not a
forecast by itself, since no population makes up all the stellar mass, but it is the one number
a reader needs to use any $F$ they prefer. Table~\ref{tab:n1} lists it for every selection the
analysis produces, in both counting conventions, and Fig.~\ref{fig:yieldF} draws
Eq.~(\ref{eq:linear}) with the literature abundances marked. The script checks that the ratio
$N/F$ is the same, to machine precision, at every grid value it was written at.

\begin{table}[h]
\centering
\caption{Yield per unit abundance: expected events over the ten-year window at $F=1$, inside
Roman's $1.47\,\mathrm{deg}^2$ footprint unless stated. Multiply by $F$ for a yield. Errors are
Monte Carlo only. The last column is the ordinary-lens run, in which every lens is drawn from
the full stellar population, so $F=1$ \emph{by definition} and its entries are yields as they
stand; its mass-precision rows are omitted because that run predates changes to the Fisher
derivatives (its detection logic is unchanged). $\eta$ is the ratio of a population's $N_1$ to
the ordinary-lens yield for the same selection: the fraction of all events that population
supplies per unit mass fraction.}
\label{tab:n1}
\small
\begin{tabular}{lrrrrr}
\toprule
 & \multicolumn{2}{c}{black holes, $N_1$} & \multicolumn{2}{c}{neutron stars, $N_1$} & ordinary \\
\cmidrule(lr){2-3}\cmidrule(lr){4-5}
selection & per object & all stars & per object & all stars & lenses \\
\midrule
either survey detects        & $5340\pm70$  & 5760 & $38\,300\pm500$ & 40\,100 & 57\,000 \\
Roman detects                & $4920\pm70$  & 5310 & $28\,100\pm400$ & 29\,600 & 41\,200 \\
Rubin detects                & $2820\pm50$  & 2940 & $19\,700\pm400$ & 20\,100 & 26\,700 \\
both detect                  & $2430\pm50$  & 2530 & $10\,200\pm300$ & 10\,400 & 12\,100 \\
Roman only                   & $2490\pm50$  & 2790 & $17\,900\pm400$ & 19\,200 & 29\,100 \\
Rubin only                   & $393\pm16$   & 417  & $9480\pm260$    & 9710    & 14\,700 \\
Roman, $\sigma(\Ml)/\Ml<1\%$  & $25.6\pm5.0$ & 25.8 & $11.5\pm2.7$    & 11.5    & --- \\
Roman, $\sigma(\Ml)/\Ml<5\%$  & $198\pm13$   & 201  & $312\pm29$      & 315     & --- \\
Roman, $\sigma(\Ml)/\Ml<10\%$ & $377\pm18$   & 388  & $708\pm47$      & 720     & --- \\
joint, $\sigma(\Ml)/\Ml<10\%$ & $520\pm21$   & 533  & $957\pm59$      & 976     & --- \\
\midrule
Rubin detects, full $59\,\mathrm{deg}^2$ & $44\,100\pm300$ & 46\,700 & $268\,000\pm2000$ & 276\,000 & 351\,000 \\
\midrule
$\eta$, Roman detects        & \multicolumn{2}{c}{0.119} & \multicolumn{2}{c}{0.683} & \\
$\eta$, Rubin detects (full scan) & \multicolumn{2}{c}{0.126} & \multicolumn{2}{c}{0.763} & \\
\bottomrule
\end{tabular}
\end{table}

\begin{figure}[h]
\centering
\includegraphics[width=\textwidth]{y3_yield_vs_F.png}
\caption{Expected events over the ten-year window inside Roman's footprint against the assumed
abundance, Eq.~(\ref{eq:linear}); the shaded bands are the Monte Carlo $1\sigma$ and are barely
wider than the lines. Dotted verticals mark the literature abundances of
Section~\ref{sec:yield}: SM26 \citet{SajadianMakler2026}; SS23 \citet{SajadianSahu2023}, their
$F_1$; O20, L20, G00 \citet{Olejak2020}, \citet{Lam2020} and \citet{Gould2000}; S22
\citet{Sweeney2022}, for neutron stars and black holes together; G00 \citet{Gould2000}. The
horizontal line is one event. Black holes over $3$--$1000\,\Msun$ as simulated; per-object
convention.}
\label{fig:yieldF}
\end{figure}

\paragraph{Reading the figure.} Roman detects black holes at a rate of $\sim\!49$ per per cent
of stellar mass, so the literature range $F_{\rm BH}=0.0045$--$0.03$ maps to $22$--$150$ events;
to expect even one black hole with its mass measured to $10\%$ needs $F_{\rm BH}\gtrsim0.002$ with
the joint fit, or $0.003$ with Roman alone. Neutron
stars are $\sim\!6\times$ more productive per unit $F$ ($28\,100$ against $4920$), which is
the $\langle\sqrt{\Ml}\rangle/\langle\Ml\rangle$ argument again: $0.85$ for neutron stars against
$0.060$ for black holes. The advantage shrinks to $\sim\!2\times$ for mass measurements, because
the larger $\tetE$ of a black hole is what makes its mass measurable.

\paragraph{The ratio $\eta$ closes the loop on Section~\ref{sec:yield}.} There we argued that
$\langle\sqrt{\Ml}\rangle/\langle\Ml\rangle$ suppresses black holes by a factor $33$ relative to
ordinary lenses at equal mass fraction, and that detection efficiency hands back about $4\times$.
The product predicts $\eta\approx4/33=0.12$; the two independent runs give $0.119$. For neutron
stars the same argument gives $0.85/1.97=0.43$ before efficiency, so $\eta=0.68$ implies a
detection efficiency $1.6\times$ that of ordinary lenses: the same direction as for black holes
and a smaller factor, as expected for lenses only $\sim\!6\times$ heavier than the typical star. For small $F$ a population supplies a fraction $\approx\eta F$
of all events: $0.12\%$ of Roman's detections are black holes at $F_{\rm BH}=0.01$.

\subsection{How to compute a yield}
\label{sec:recipe}

Everything above is computed by two scripts from the per-event output table, and none of it
requires a new simulation unless the population itself changes. The procedure, with the
choices it makes explicit:

\begin{enumerate}
\item \textbf{Simulate the population.} Run the simulator with every lens drawn from the population
  under test, e.g.\ \verb|./roman --population bh --events 300 --lenses 50 --stride-roman 5|,
  and archive its output tree as a run directory. The table must keep the \emph{undetected}
  draws: $\langle\Ml\rangle$ is taken over all draws of a lens component, and the detected ones
  are biased heavy.
\item \textbf{Attach a rate to each draw.} For draw $i$ on sightline $k$,
  \begin{equation*}
  y_i \;=\; T\;\frac{w_{{\rm area},k}\,N_{\star,k}}{n_{{\rm sim},k}}\;
            \frac{2u_{0,\max}\,\kappa}{\langle\Ml\rangle}\;Z(D_{s,i})\,\sqrt{M_{{\rm L},i}}\,v_{t,i},
  \end{equation*}
  with $T$ the window over which $t_0$ is drawn ($10\,\mathrm{yr}-4\,\mathrm{d}$), $w_{\rm area}$ the
  sightline's sky area (which undoes the finer grid inside Roman's footprint), $N_\star$ its
  source-star count, $n_{\rm sim}$ the number of draws it made, $u_{0,\max}=3$, and
  $\kappa=\sqrt{4G/c^2}$. $y_i$ is the expected number of events draw $i$ stands for, at $F=1$,
  per resolved object.
\item \textbf{Choose the convention.} Keep $y_i$ to count per resolved object, as the Sajadian
  papers do, or divide by the probability the simulator kept the draw,
  $P = 1-(1-f_{b,r}\,[\text{Rubin could see it}])(1-f_{b,\rm F146}\,[\text{Roman could}])$, to count
  on every star.
\item \textbf{Select and sum.} $N_1=\sum_i y_i\,s_i$, where $s_i\in\{0,1\}$ is whatever selection is
  wanted: detected by Roman, by both, mass to $10\%$, inside the footprint. The Monte Carlo error
  is $(\sum_i y_i^2 s_i)^{1/2}$.
\item \textbf{Multiply by $F$}, Eq.~(\ref{eq:linear}).
\item \textbf{Check the normalisation.} Before trusting $N_1$, rebuild the optical depth from the
  same constants, $\tau=(\pi/2u_{0,\max})\langle\gamma\,\tE\rangle$, and compare it with the
  optical depth the simulator computes independently for each sightline. A factor error in any
  constant, unit or in $Z$ shows up here; ours agree to $0.1\%$.
\end{enumerate}

Steps 2--6 are \texttt{analysis/y1\_absolute\_yield.py}, which writes the tables of this section
and a CSV; \texttt{analysis/y3\_yield\_vs\_F.py} reads that CSV, extracts $N_1$, and draws
Fig.~\ref{fig:yieldF}:

\begin{verbatim}
.roman/bin/python analysis/y1_absolute_yield.py \
    --run bh=runs/prod_bh_20260917 --run ns=runs/prod_ns_20260917 \
    --run bulge=../roman_runs/2026-09-06_v3_h7 \
    --F bh=0.005,0.01,0.03 --F ns=0.005,0.01,0.03,0.06 \
    -o figures/yield_prefix_20260922
.roman/bin/python analysis/y3_yield_vs_F.py \
    figures/yield_prefix_20260922/y1_yields.csv -o figures/yield_prefix_20260922/y3
\end{verbatim}

\paragraph{Other windows, areas and mass functions.}
\begin{itemize}
\item \emph{A different window or area.} Both are selections in step~4, on $t_0$ or on sightline.
  Do not rescale by $T$: the detection efficiency depends on where $t_0$ falls relative to the
  seasons, so a yield is not proportional to the window length. Per-$\mathrm{deg}^2$ figures do
  follow by division.
\item \emph{A different mass function}, without re-running. Re-weight each draw by the ratio of
  the new to the sampled mass density, $p_{\rm new}(M_i)/p_{\rm old}(M_i)$, and replace
  $\langle\Ml\rangle$ by the new mean. The $3$--$50\,\Msun$ rows of Table~\ref{tab:yieldbh} are the
  simplest case: a narrower log-uniform range keeps only the draws inside it and scales them by
  $[\ln(1000/3)/\ln(50/3)]\times\langle\Ml\rangle_{3-1000}/\langle\Ml\rangle_{3-50}$. This is
  exact only where the sampled function covers the new one, and it costs precision in proportion
  to how unevenly the new weights fall (the Kish $\Neff$ of the re-weighted sample says how much).
\item \emph{Black holes and neutron stars together.} The two runs are separate simulations, so
  their yields add: $N = F_{\rm BH}N_{1,\rm BH} + F_{\rm NS}N_{1,\rm NS}$. What must not be done is
  to pool their draws in one weighted statistic, because each weight carries the $\sqrt{\Ml}$ of
  its own mass function.
\end{itemize}

\subsection{Comparison with \texorpdfstring{\citet{SajadianSahu2023}}{Sajadian \& Sahu (2023)}}

Their method is not ours and the two must be aligned before any number is compared. They take
Roman's \emph{total} event count $N_{e,\rm tot}=27{,}000$ from \citet{Penny2019} as given, and
multiply it by two factors: $F_1 = 0.019$, the fraction of stellar mass in black holes, obtained
from the Besan\c{c}on disc initial mass function by converting every star born above $20\,\Msun$
into a remnant; and $F_2$, the ratio of $\langle\epsilon(\tE)/\tE\rangle$ for black holes to the
same quantity for all lenses, which corrects for the fact that black-hole events last far longer
and so occur at a lower rate per unit optical depth. Two points follow immediately. First,
\emph{their $F_1$ is the same kind of quantity as our $F$}, so the honest comparison is at
$F = 0.019$, not at the Sajadian \& Makler value of $0.005$. Second, their $56$--$77$ is the
\emph{uniform} mass-function result; the row that matches our log-uniform function is
$\mathrm{d}N/\mathrm{d}M \propto M^{-1}$, for which $F_2 = 0.16$--$0.17$ and
$N_{e,\rm BHs} \approx 86$.

\paragraph{Detections.} Matching mass function, abundance, area and season count, we exceed them
by a factor $\sim\!2.3$:

\begin{center}
\begin{tabular}{lrrr}
\toprule
 & detections & deg$^2$ $\times$ seasons & per deg$^2$ per season \\
\midrule
\citet{SajadianSahu2023}, $\alpha=1$, $F_1=0.019$ & 86  & $1.96\times6$  & 7.2 \\
this work, $3$--$50\,\Msun$, $F=0.019$           & 245 & $1.47\times10$ & 16.7 \\
\bottomrule
\end{tabular}
\end{center}

Part of the excess is a difference in the detection cut: they require
$\Delta\chi^2 > 800$ \emph{and} at least three points $4\sigma_m$ above baseline, against our
$\Delta\chi^2 \ge 500$, so their bar is the higher one. The rest is most likely the survey design
itself, and that is a point about vintage rather than about either calculation. \citet{Penny2019} describe the Cycle~7 architecture --
$1.96\,\mathrm{deg}^2$ observed in $6\times72$ days -- and Sajadian \& Sahu adopt that observing
strategy wholesale. Its defining feature, for them, is a $\sim\!2.3$-year gap between the first
three and the last three seasons, which they identify as the main thing suppressing black-hole
characterisation, and which they patch by hand with one hour of observing every ten days. That
patch alone raises their $F_2$ from $0.11$ to $0.15$, a $36\%$ gain in detections from
approximately one day of extra telescope time. The design this simulator follows \citep{ROTAC2025}
\emph{fills that gap as a matter of course}, with four low-cadence seasons between the six
high-cadence ones. We therefore expect to detect more long events than they do, which is exactly
what the factor $2.3$ is. Neither their number nor Penny's should be read as a forecast for the
survey as finally designed; both predate it.

A caveat that runs the other way: their $F_1$ equates the progenitor mass with the remnant mass,
which overstates the black-hole mass fraction, so $0.019$ is an upper bound on the $F$ they
effectively assume.

\paragraph{Characterisation.} Our characterised counts fall below theirs, and locating the
deficit requires care with two things the obvious comparison gets wrong. Their percentages are
fractions of \emph{detectable} events, whereas Table~\ref{tab:tenpercent} of this report is a
fraction of \emph{characterised} events, a smaller denominator; and their mass range is
$2$--$50\,\Msun$ against our $3$--$1000$. Redoing the comparison on their denominator and on the
$3$--$50\,\Msun$ subset (\texttt{analysis/y2\_ss23\_compare.py}), at the $10\%$ threshold:

\begin{center}
\begin{tabular}{lrrrr}
\toprule
 & Sajadian \& Sahu [\%] & this work [\%] & ratio & (full $3$--$1000$) \\
\midrule
$\sigma(\tetE)/\tetE < 10\%$ & 99.2 & 68.6 & 0.69 & 84.0 \\
$\sigma(\tE)/\tE < 10\%$     & 67.0 & 21.7 & 0.32 & 20.1 \\
$\sigma(\piE)/\piE < 10\%$   & 30.0 &  7.4 & 0.25 &  8.3 \\
$\sigma(\Ml)/\Ml < 10\%$     & 29.8 &  6.2 & 0.21 &  7.7 \\
\bottomrule
\end{tabular}
\end{center}

Every parameter falls short, but not equally, and the ordering is the informative part.
$\tetE$ --- the purely astrometric quantity, and the one our $1.1$\,mas per-exposure floor
controls --- is the \emph{least} deficient at $0.69$, while $\piE$ and the lens mass that depends
on it sit at $0.25$ and $0.21$. The gap is therefore concentrated on the parallax rather than on
the astrometry, and the difference in assumed astrometric precision --- they use
magnitude-dependent jitter simulations (S.\ Calchi Novati, private communication), we use a
$1.1$\,mas floor rising to $10$\,mas at $F146 = 23.5$ --- is unlikely to be the principal cause.

The last column also disposes of the explanation we first reached for. If our deficit in $\tE$ and
$\piE$ came from the $3$--$1000\,\Msun$ range producing events too long for Roman's $\sim\!70$-day
seasons, restricting to $3$--$50\,\Msun$ should have repaired it. It does not: $\tE$ moves from
$20.1$ to $21.7\%$ and $\piE$ from $8.3$ to $7.4\%$, both within noise. What the restriction does
change is $\tetE$, which \emph{falls} from $84.0$ to $68.6\%$ --- as it must, since
$\tetE \propto \sqrt{\Ml}$ and the subset is lighter. The mass function is not the cause.

Our Roman-partition photometric matrix has the same seven free parameters as theirs
($t_0,u_0,\tE,\xi,f_b,m_{\rm base},\piE$), so this is not a difference of parametrisation there.
The astrometric matrices do differ: theirs has three parameters ($\tetE,\mu_{s,1},\mu_{s,2}$) and
ours four, because we fit $\piE$ astrometrically as well. That should help $\piE$ rather than hurt
it, so it does not explain the deficit either. We record the parallax shortfall as an open
question rather than claim a cause for it.

\paragraph{These yields precede the extinction correction.} Every number in this report was
produced with an inverted \citet{Cardelli1989} reddening law, which evaluated the
extinction at $1/\lambda$ rather than at $\lambda$ (Deviation 53). The error runs in
\emph{opposite directions for the two telescopes}: in the optical it gave $A_u/A_V = 0.07$ against
the correct $1.694$, making Rubin's sources far too bright, while in the infrared it gave
$A_{F146}/A_V = 0.751$ against the correct $0.197$, making Roman's sources too faint. The Rubin
yields above are therefore biased \emph{high} and the Roman yields \emph{low}, and the report's
finding that Rubin dominates the detection count is biased in its own favour. A controlled stub
run across the fix moved Roman-only detections from $31$ to $41$ while leaving Rubin-only
unchanged at $42$. The corrected production runs are reported separately.

\section{Caveats}

\begin{itemize}
\item \textbf{The populations are run separately and must not be pooled.} The event-rate weight
      contains $\sqrt{\Ml}$, which is valid only for the mass function actually sampled. Any
      combined black-hole plus neutron-star statement requires an assumed relative abundance,
      applied after the fact.
\item \textbf{The log-uniform black-hole mass function is an assumption, not a measurement.} It
      is taken as the truth here rather than as an importance-sampling device.
\item \textbf{Absolute yields scale linearly with an assumed abundance $F$, which is uncertain by
      an order of magnitude.} Section~\ref{sec:yield} tabulates a grid rather than a single number
      for exactly this reason; the literature spans $F_{\rm BH} = 0.005$ to $0.03$.
\item \textbf{The extinction law was inverted when these runs were made} (Deviation 53). It biases
      Rubin's yields high and Roman's low, in opposite directions, so the survey-share numbers in
      Table~\ref{tab:synergy} favour Rubin more than they should. Corrected runs are reported
      separately.
\item \textbf{The model's latitude gradient is too shallow.} Against OGLE-IV it reproduces the
      event rate near $b=-1.5^\circ$ but overproduces by $\sim\!40\%$ at $b=-5^\circ$, so
      full-scan yields carry a systematic of a few tens of per cent. Footprint yields are much
      less affected.
\item \textbf{The footprint sample is small.} Every Roman-versus-joint comparison rests on
      $\Neff \approx 5400$, against $\Neff \approx 1.4$--$3.0\times10^5$ for the full scan.
\item \textbf{$\sigma_{\rm joint} \le \sigma_{\rm single}$ is violated on six events} by a factor
      $1.0003$, consistent with round-off in the $9\times9$ matrix inversion on ill-conditioned
      events rather than with an error in the information sum.
\item \textbf{Rubin's astrometric error model is the weakest input.} The delivered per-visit
      floor is taken as $10$\,mas; published estimates range from $3$ to $7$\,mas, which would
      make Rubin better than assumed here and would raise its $\tetE$ column specifically.
\end{itemize}

%\section*{Reproducing this report}
%
%\begin{verbatim}
%.roman/bin/python analysis/p6_synergy_resolution.py \
%    --run bh=runs/prod_bh_20260917 --run ns=runs/prod_ns_20260917 \
%    -o figures/prod_20260918/p6
%.roman/bin/python analysis/p7_forecast_figures.py \
%    --run bh=runs/prod_bh_20260917 --run ns=runs/prod_ns_20260917 \
%    -o figures/prod_20260918/p7
%cd Report/populations && latexmk -pdf populations_report.tex
%\end{verbatim}

\bibliographystyle{plainnat}
\bibliography{../refs}

\end{document}
